\section{AHSME 1972}
        \begin{enumerate}
        	\item (\textit{AHSME 1972}) The lengths in inches of the three sides of each of four triangles $I, II, III$, and $IV$ are as follows: 


 $\begin{array}{rlrl} \hbox{I}& 3,\ 4,\ \hbox{and}\ 5\qquad & \hbox{III}& 7,\ 24,\ \hbox{and}\ 25\\  \hbox{II}& 4,\ 7\frac{1}{2},\ \hbox{and}\ 8\frac{1}{2}\qquad & \hbox{IV}& 3\frac{1}{2},\ 4\frac{1}{2},\ \hbox{and}\ 5\frac{1}{2}.\end{array}$


 Of these four given triangles, the only right triangles are 


 $\text{(A)  I and II}\qquad  \text{(B)  I and III}\qquad \text{(C)  I and IV} \quad \\ \text{(D)  I, II, and III} \qquad \text{(E)  I, II, and IV }$

 \textbf{Answer: }D



\textbf{Solution}

Let's test out the four triangles:

 I: $3^2 + 4^2 = 5^2$ (true)

 II: $4^2 + 7.5^2 = 8.5^2$ (true)

 III: $7^2 + 24^2 = 25^2$

 The only answer choice with I, II, and III in the answer is $\fbox{D} \rightarrow$ I, II, and III.

 


	\item (\textit{AHSME 1972}) If a dealer could get his goods for $8$\% less while keeping his selling price fixed, 
his profit, based on cost, would be increased to $(x+10)$\% from his present profit of $x$\%, which is


 $\textbf{(A) }12\%\qquad \textbf{(B) }15\%\qquad \textbf{(C) }30\%\qquad \textbf{(D) }50\%\qquad  \textbf{(E) }75\%$

 \textbf{Answer: }B



\textbf{Solution}

If $c$ is the cost, then $0.92c$ is the reduced cost. Selling price is equal to cost plus profit:

 

 \[c(1+0.01x) = 0.92c(1 + 0.01(x+10))\]

 Solving for $x$ yields $x = 15.$ The answer is $\textbf{(B)}.$

 -edited by coolmath34

 


	\item (\textit{AHSME 1972}) If $x=\dfrac{1-i\sqrt{3}}{2}$ where $i=\sqrt{-1}$, then $\dfrac{1}{x^2-x}$ is equal to


 $\textbf{(A) }-2\qquad \textbf{(B) }-1\qquad \textbf{(C) }1+i\sqrt{3}\qquad \textbf{(D) }1\qquad  \textbf{(E) }2$

 \textbf{Answer: }B



\textbf{Solution}

Using DeMoivre's theorem, we can calculate $x^2=\frac{1+i\sqrt{3}}{2}$ The denominator is therefore $-1$ which makes the answer 
 \[\boxed{\textbf{(B) }-1}.\] ~lopkiloinm

 


	\item (\textit{AHSME 1972}) The number of solutions to $\{1,~2\}\subseteq~X~\subseteq~\{1,~2,~3,~4,~5\}$, where $X$ is a subset of $\{1,~2,~3,~4,~5\}$ is


 $\textbf{(A) }2\qquad \textbf{(B) }4\qquad \textbf{(C) }6\qquad \textbf{(D) }8\qquad  \textbf{(E) }\text{None of these}$

 \textbf{Answer: }D



\textbf{Solution}

$X$ has to contain $\{1,~2\}$, so only $\{3,~4,~5\}$ matters. There are two choices for the elements; the element is either in $X$ or outside of $X$. With this combinatorics in mind, the answer is simply $2^3=\boxed{\textbf{(D) }8}.$ ~lopkiloinm

 


	\item (\textit{AHSME 1972}) From among $2^{1/2}, 3^{1/3}, 8^{1/8}, 9^{1/9}$ those which have the greatest and the next to the greatest values, in that order, are 


 $\textbf{(A) } 3^{1/3},\ 2^{1/2}\quad  \textbf{(B) } 3^{1/3},\ 8^{1/8}\quad  \textbf{(C) } 3^{1/3},\ 9^{1/9}\quad  \textbf{(D) } 8^{1/8},\ 9^{1/9}\quad \\  \text{(E) None of these}$

 \textbf{Answer: }A



\textbf{Solution}

$8^{1/8}$ and $9^{1/9}$ are obviously too small. We must then compare $3^{1/3}$ with $2^{1/2}$. Raising both to the power of $6$ gives $9$ and $8$ respectively. Our answer is therefore $\boxed{\textbf{(A) } 3^{1/3},\ 2^{1/2}}.$ ~lopkiloinm

 


	\item (\textit{AHSME 1972}) If $3^{2x}+9=10\left(3^{x}\right)$, then the value of $(x^2+1)$ is 


 $\textbf{(A) }1\text{ only}\qquad \textbf{(B) }5\text{ only}\qquad \textbf{(C) }1\text{ or }5\qquad \textbf{(D) }2\qquad \textbf{(E) }10$

 \textbf{Answer: }C



\textbf{Solution}

We can rearrange to get $3^{2x}-10\left(3^{x}\right)+9=0$. Factor or use the quadratic formula to find that $3^{x}= 9\text{ or }1$, which makes $x = 2\text{ or }0$ and our answer 
 \[\boxed{\textbf{(C) }1\text{ or }5}.\]



	\item (\textit{AHSME 1972}) If $yz:zx:xy=1:2:3$, then $\dfrac{x}{yz}:\dfrac{y}{zx}$ is equal to


 $\textbf{(A) }3:2\qquad \textbf{(B) }1:2\qquad \textbf{(C) }1:4\qquad \textbf{(D) }2:1\qquad  \textbf{(E) }4:1$

 \textbf{Answer: }E



\textbf{Solution}

We want to find
 \[\dfrac{x}{yz}:\dfrac{y}{zx} = \dfrac{x}{yz} \cdot \dfrac{zx}{y} = \dfrac{x^2}{y^2}.\]

 We are given $yz:zx=1:2,$ and dividing both sides by $z$ gives $y:x = 1:2.$ Therefore,
 \[\dfrac{x^2}{y^2} = \dfrac{1}{\dfrac{1}{2^2}} = 4.\]

 The answer is $\textbf{(E)}.$

 -edited by coolmath34

 


	\item (\textit{AHSME 1972}) If $|x-\log y|=x+\log y$ where $x$ and $\log y$ are real, then


 $\textbf{(A) }x=0\qquad \textbf{(B) }y=1\qquad \textbf{(C) }x=0\text{ and }y=1\qquad\\ \textbf{(D) }x(y-1)=0\qquad  \textbf{(E) }\text{None of these}$

 \textbf{Answer: }D



\textbf{Solution}

We have two cases: $x$ is positive and $x$ is negative.
 \[\text{Positive} \qquad x - \log y = x + \log y \qquad \rightarrow \qquad 2\log y = 0 \rightarrow y = 1\]
 \[\text{Negative} \qquad -x + \log y = x + \log y \qquad \rightarrow \qquad 2x = 0 \rightarrow x = 0\]

 We can write $x(y-1)=0$ to show that $x=0, y=1,$ or both.

 The answer is $\textbf{(D)}.$

 -edited by coolmath34

 


	\item (\textit{AHSME 1972}) Ann and Sue bought identical boxes of stationery. Ann used hers to write $1$-sheet letters and Sue used hers to write $3$-sheet letters. 
Ann used all the envelopes and had $50$ sheets of paper left, while Sue used all of the sheets of paper and had $50$ envelopes left. 
The number of sheets of paper in each box was 


 $\textbf{(A) } 150 \qquad \textbf{(B) } 125 \qquad \textbf{(C) } 120 \qquad \textbf{(D) } 100 \qquad     \textbf{(E) } 80 \qquad$

 \textbf{Answer: }A



\textbf{Solution}

Let $S$ be the number of sheets of paper and $E$ be the number of envelopes. We can write two equations:
 \[S - E = 50\]
 \[E - S/3 = 50\]

 Solving yields $(E, S) = (100, 150).$

 The answer is $\textbf{(A)}.$

 -edited by coolmath34

 


	\item (\textit{AHSME 1972}) For $x$ real, the inequality $1\le |x-2|\le 7$ is equivalent to


 $\textbf{(A) }x\le 1\text{ or }x\ge 3\qquad \textbf{(B) }1\le x\le 3\qquad \textbf{(C) }-5\le x\le 9\qquad \\ \textbf{(D) }-5\le x\le 1\text{ or }3\le x\le 9\qquad  \textbf{(E) }-6\le x\le 1\text{ or }3\le x\le 10$

 \textbf{Answer: }D



\textbf{Solution}

We can split the inequality into two smaller inequalities and solve them individually.
 \[|x-2| \ge 1 \quad \rightarrow \quad x \ge 3 \quad \text{and} \quad x \le 1\]
 \[|x-2| \le 7 \quad \rightarrow \quad x \le 9 \quad \text{and} \quad x \ge -5\]

 Combining these inequalities, we get $x \in [-5, 1] \cup [3, 9].$

 The answer is $\textbf{(D)}.$

 -edited by coolmath34

 


	\item (\textit{AHSME 1972}) The value(s) of $y$ for which the following pair of equations $x^2+y^2-16=0\text{ and }x^2-3y+12=0$ may have a real common solution, are


 $\textbf{(A) }4\text{ only}\qquad \textbf{(B) }-7,~4\qquad \textbf{(C) }0,~4\qquad \textbf{(D) }\text{no }y\qquad  \textbf{(E) }\text{all }y$

 \textbf{Answer: }A



\textbf{Solution (Quick)}

Checking the answer choices, we see that only 4 is the viable choice. Therefore, the answer is $\boxed{A}$



	\item (\textit{AHSME 1972}) The number of cubic feet in the volume of a cube is the same as the number of square inches in its surface area. The length of the edge expressed as a number of feet is


 $\textbf{(A) }6\qquad \textbf{(B) }864\qquad \textbf{(C) }1728\qquad \textbf{(D) }6\times 1728\qquad  \textbf{(E) }2304$

 \textbf{Answer: }B



\textbf{Solution}

If the side length of a cube is $s$ feet the volume is $s^3$ cubic feet and the surface area is $6(12s^2)$ square inches. The volume and surface are numerically equal, so we can write
 \[s^3 = 6(12s^2)\]Solving yields $s=864.$

 The answer is $\textbf{(B)}.$



	\item (\textit{AHSME 1972}) 
                    \begin{figure}[H]
                        \centering
                        \includegraphics[width=0.5\textwidth]{Problem/AHSME_I_1972_13_Problem_Figure_1.png}
                    \end{figure}
                    


 Inside square $ABCD$ (See figure) with sides of length $12$ inches, segment $AE$ is drawn where $E$ is the point on $DC$ which is $5$ inches from $D$. 
The perpendicular bisector of $AE$ is drawn and intersects $AE, AD$, and $BC$ at points $M, P$, and $Q$ respectively. The ratio of segment $PM$ to $MQ$ is


 $\textbf{(A) }5:12\qquad \textbf{(B) }5:13\qquad \textbf{(C) }5:19\qquad \textbf{(D) }1:4\qquad  \textbf{(E) }5:21$

 \textbf{Answer: }C



\textbf{Solution}

Let the line passing through $M$ parallel to $AB$ intersect $AD$ and $BC$ and $S$ and $T$ respectively. Since $M$ is the midpoint of $AE$, $SM=\frac{5}{2}$ and $TM=12-\frac{5}{2}=\frac{19}{2}$. Since $\triangle PSM\sim \triangle QTM$, $PM:MQ=SM:MT=5:19$, hence our answer is $\fbox{C}$.

 


	\item (\textit{AHSME 1972}) A triangle has angles of $30^\circ$ and $45^\circ$. If the side opposite the $45^\circ$ angle has length $8$, then the side opposite the $30^\circ$ angle has length


 $\textbf{(A) }4\qquad \textbf{(B) }4\sqrt{2}\qquad \textbf{(C) }4\sqrt{3}\qquad \textbf{(D) }4\sqrt{6}\qquad  \textbf{(E) }6$

 \textbf{Answer: }B



\textbf{Solution}

This triangle can be split into smaller 30-60-90 and 45-45-90 triangles. The side opposite the $45^\circ$ angle has length $8,$ so the 30-60-90 triangle has sides $4, 4\sqrt3,$ and $8.$

 One of the legs of the 45-45-90 triangles is $4,$ so the hypotenuse is $4\sqrt2.$ This is also the side opposite the $30^\circ$ angle, so the answer is $\textbf{(B)}.$

 -edited by coolmath34

 



\textbf{Solution 2}

Use \href{https://artofproblemsolving.com/wiki/index.php?title=Law\_of\_Sines}{Law of Sines} to get $\frac{8}{\frac{\sqrt2}{2}}=\frac{x}{\frac{1}{2}}$. The answer is $\boxed{\textbf{(B) }4\sqrt{2}}$.

 -aopspandy

 


	\item (\textit{AHSME 1972}) A contractor estimated that one of his two bricklayers would take $9$ hours to build a certain wall and the other $10$ hours. 
However, he knew from experience that when they worked together, their combined output fell by $10$ bricks per hour. 
Being in a hurry, he put both men on the job and found that it took exactly 5 hours to build the wall. The number of bricks in the wall was


 $\textbf{(A) }500\qquad \textbf{(B) }550\qquad \textbf{(C) }900\qquad \textbf{(D) }950\qquad  \textbf{(E) }960$

 \textbf{Answer: }C



\textbf{Solution}

Every hour, the contractors lay $x/9$ and $x/10$ bricks, assuming a total of $x$ bricks in the wall. Together, $x/5$ bricks are laid per hour:
 \[x = 5(x/9 + x/10 - 10)\]There are a total of $x=900$ bricks, and the answer is $\textbf{(C)} \quad 900$.

 -edited by coolmath34

 


	\item (\textit{AHSME 1972}) There are two positive numbers that may be inserted between $3$ and $9$ such that the first three are in geometric progression 
while the last three are in arithmetic progression. The sum of those two positive numbers is


 $\textbf{(A) }13\textstyle\frac{1}{2}\qquad \textbf{(B) }11\frac{1}{4}\qquad \textbf{(C) }10\frac{1}{2}\qquad \textbf{(D) }10\qquad \textbf{(E) }9\frac{1}{2}$

 \textbf{Answer: }B



\textbf{Solution 1}

Let $a$ be first and $b$ be second. We can then get equations based on our knowledge: $b-a = 9-b$ and $b/a = a/3$. We then get $b = (9+a)/2$ from our first equation and we substitute that into the second to get $\frac{9+a}{2a} = a/3$ which simplifies to $2a^2-3a-27=0$ which be $(2a-9)(a+3) = 0$ and $a=9/2$. Then $b=27/4$. Their sum is $\boxed{\textbf{(B) }11\frac{1}{4}}$. ~lopkiloinm

 



\textbf{Solution 2}

Like in Solution 1, let $a$ be the first number and $b$ be the second. From the problem, we have that

 
$a = 3r$ $\textbf{(1)}$
$b = 3r^2$ $\textbf{(2)}$
$9 = a + 2d$ $\textbf{(3)}$

 
for some common difference $d$ and common ratio $r$. We can use these to obtain the following:

 
$3r + d = 3r^2$ (from $\textbf{(1)}$, $\textbf{(2)}$, and the definition of an arithmetic sequence) $\textbf{(4)}$
$9 = 3r + 2d$ (from $\textbf{(1)}$ and $\textbf{(3)}$) $\textbf{(5)}$

 
Solving $\textbf{(5)}$ for $d$ and substituting into $\textbf{(4)}$, we find that $3r + \frac{9 - 3r}{2} = 3r^2$ or $6r + (9 - 3r) = 6r^2$. Moving every term to the right side and factoring, we get that $3(2r - 3)(r + 1) = 0$. Because $a$ and $b$ must be positive, the common ratio cannot be $-1$, so we ascertain that the common ratio is $3/2$. From here, we find that $a = 3 \cdot \frac{3}{2} = \frac{9}{2}$, and $b = \frac{9}{2} \cdot \frac{3}{2} = 27/4$. Then their sum is $\boxed{\textbf{(B) }11\frac{1}{4}}$.

 


	\item (\textit{AHSME 1972}) A piece of string is cut in two at a point selected at random. The probability that the longer piece is at least x times as large as the shorter piece is


 $\textbf{(A) }\frac{1}{2}\qquad \textbf{(B) }\frac{2}{x}\qquad \textbf{(C) }\frac{1}{x+1}\qquad \textbf{(D) }\frac{1}{x}\qquad \textbf{(E) }\frac{2}{x+1}$

 \textbf{Answer: }E



\textbf{Solution}

The string could be represented by a line plot segment on the interval $(0,x+1)$. The interval $(0,1)$ is a solution as well as $(x,x+1)$. The probability is therefore $\boxed{\textbf{(E) }\frac{2}{x+1}}.$ ~lopkiloinm

 


	\item (\textit{AHSME 1972}) Let $ABCD$ be a trapezoid with the measure of base $AB$ twice that of base $DC$, and let $E$ be the point of intersection of the diagonals. 
If the measure of diagonal $AC$ is $11$, then that of segment $EC$ is equal to


 $\textbf{(A) }3\textstyle\frac{2}{3}\qquad \textbf{(B) }3\frac{3}{4}\qquad \textbf{(C) }4\qquad \textbf{(D) }3\frac{1}{2}\qquad \textbf{(E) }3$

 \textbf{Answer: }A



\textbf{Solution}

We begin with a diagram:

 
                    \begin{figure}[H]
                        \centering
                        \includegraphics[width=0.5\textwidth]{Solution/AHSME_I_1972_18_Solution_Figure_1.png}
                    \end{figure}
                    

 The bases of a trapezoid are parallel by definition, so $\angle EDC$ and $\angle EBA$ are alternate interior angles, and therefore equal. We have the same setup with $\angle ECD$ and $\angle EAB$, meaning that $\triangle ABE \sim \triangle CDE$ by AA Similarity. We could've also used the fact that $\angle BEA$ and $\angle DEC$ are vertical angles.

 With this information, we can setup a ratio of corresponding sides:
 \[\frac{AB}{CD} = \frac{AE}{CE} \implies \frac{2CD}{CD} = \frac{11 - CE}{CE}.\]And simplify from there:
 \begin{align*} \frac{2CD}{CD} &= \frac{11 - CE}{CE} \\ 2 &= \frac{11 - CE}{CE} \\ 2CE &= 11 - CE \\ 3CE &= 11 \\ CE &= \frac{11}{3} = 3 \frac{2}{3}. \end{align*}
Therefore, our answer is $\boxed{\textbf{(A) }3\textstyle\frac{2}{3}.}$



	\item (\textit{AHSME 1972}) The sum of the first $n$ terms of the sequence $1,~(1+2),~(1+2+2^2),~\dots ~(1+2+2^2+\dots +2^{n-1})$ in terms of $n$ is


 $\textbf{(A) }2^n\qquad \textbf{(B) }2^n-n\qquad \textbf{(C) }2^{n+1}-n\qquad \textbf{(D) }2^{n+1}-n-2\qquad  \textbf{(E) }n\cdot 2^n$

 \textbf{Answer: }D



\textbf{Solution}

We begin by evaluating the sum of the first $n$ terms of the sequence with $1$, $2$, and $3$. We find that the sum of the first term of the sequence is $1$, the sum of the first two terms is $4$, and the sum of the first three terms $11$. From here, we can simply guess and check the answer choices and will find that only $\boxed{D}$ works.

 


	\item (\textit{AHSME 1972}) If $\tan x=\dfrac{2ab}{a^2-b^2}$ where $a>b>0$ and $0^\circ <x<90^\circ$, then $\sin x$ is equal to


 $\textbf{(A) }\frac{a}{b}\qquad \textbf{(B) }\frac{b}{a}\qquad \textbf{(C) }\frac{\sqrt{a^2-b^2}}{2a}\qquad \textbf{(D) }\frac{\sqrt{a^2-b^2}}{2ab}\qquad \textbf{(E) }\frac{2ab}{a^2+b^2}$

 \textbf{Answer: }E



\textbf{Solution}

We start by letting $\tan x = \frac{\sin x}{\cos x}$ so that our equation is now: 
 \[\frac{\sin x}{\cos x} = \frac{2ab}{a^2-b^2}\] Multiplying through and rearranging gives us the equation: 
 \[\cos x = \frac{a^2-b^2}{2ab} * \sin x\] We now apply the Pythagorean identity $\sin ^2 x + \cos ^2 x =1$, using our substitution: 
 \[\left(\frac{a^2-b^2}{2ab} * \sin x \right)^2 + \sin ^2 x =1\] We can isolate $\sin x$ without worrying about division by $0$ since $a \neq b \neq 0$ our final answer is $(E) \frac{2ab}{a^2+b^2}$



	\item (\textit{AHSME 1972}) 
                    \begin{figure}[H]
                        \centering
                        \includegraphics[width=0.5\textwidth]{Problem/AHSME_I_1972_21_Problem_Figure_1.png}
                    \end{figure}
                    


 If the sum of the measures in degrees of angles $A,~B,~C,~D,~E$ and $F$ in the figure above is $90n$, then $n$ is equal to


 $\textbf{(A) }2\qquad \textbf{(B) }3\qquad \textbf{(C) }4\qquad \textbf{(D) }5\qquad \textbf{(E) }6$

 \textbf{Answer: }C



\textbf{Solution}

Let $P$ be the intersection of $AC$ and $BD$. 

 Since $\angle APD = \angle BPC$ are vertical angles, the angle sums of $\triangle PAD$ and $\triangle PBC$ give:
 \[\angle A + \angle D = 180^\circ - \angle APD\]
 \[\angle B + \angle C = 180^\circ - \angle BPC\]

 Adding these equations yields:
 \[\angle A + \angle B + \angle C + \angle D = 360^\circ - 2\angle BPC\]

 Consider quadrilateral $PFE C$. Its interior angle sum gives $\angle E + \angle F + \angle BPC + \text{adj angles} = 360^\circ$, which simplifies the total sum of all six angles to:
 \[\angle A + \angle B + \angle C + \angle D + \angle E + \angle F = 360^\circ\]

 We are given that the sum equals $90n^\circ$:
 \[90n = 360 \implies n = 4\]

 Thus, $n$ is equal to $4$, which corresponds to $\mathbf{(C)}$.

 


	\item (\textit{AHSME 1972}) If $a\pm bi~(b\neq 0)$ are imaginary roots of the equation $x^3+qx+r=0$ where $a,~b,~q$, and $r$ are real numbers, then $q$ in terms of $a$ and $b$ is


 $\textbf{(A) }a^2+b^2\qquad \textbf{(B) }2a^2-b^2\qquad \textbf{(C) }b^2-a^2\qquad \textbf{(D) }b^2-2a^2\qquad  \textbf{(E) }b^2-3a^2$

 \textbf{Answer: }E



\textbf{Solution}

We have that $x^3+qx+r$ is a cubic, indicating that it has 3 roots. Two of these roots is given to you, namely $a \pm bi$, where $i = \sqrt{-1}$ so call the third root $\omega$.

 We then have the polynomial $x^3+qx+r$ to be equivalent in the form
 \[(x - (a + bi))(x - (a - bi))(x - \omega)\]Which can be simplified in the associative property into
 \[((x - a) + bi)((x - a) - bi)(x - \omega)\]

 In which if we set $\alpha$ as \( x-a \) and $\beta$ as $bi$ we can use difference of squares to get
 \[\implies (\alpha + \beta))(\alpha - \beta))(x - \omega)\]
 \[\implies (\alpha^2 - \beta^2)(x-\omega)\]
 \[\text{substitution of } \alpha \text{ and } \beta\]
 \[\implies((x-a)^2-(bi)^2)(x-\omega)\]
 \[\implies(x^2-2ax+a^2+b^2)(x-\omega)\]Expanding gives us
 \[\implies x^2(x-\omega)-2ax(x-\omega) +a^2(x-\omega)+b^2(x-\omega)\]
 \[\text{in which simplification and rearranging gives us}\]
 \[x^3-x^2\omega-2ax^2+2ax\omega+a^2x+b^2x-a^2\omega-b^2\omega\]Notice how there exists no \( x^2 \) term, indicating that the coefficients of the \( x^2 \) term must equal 0. Consequently, we get 
 \[(-\omega-2a)x^2=0x^2\]
 \[-\omega-2a=0\]
 \[-2a = \omega\]Additionally, notice how \( q \) is the coefficient of the \( x \) term, so we have
 \[2ax\omega+a^2x+b^2x = qx\]
 \[(2a\omega +a^2+b^2)x=qx\]
 \[2a\omega + a^2 + b^2 = q\]
 \[\text{using definition of } \omega\]
 \[2a(-2a)+a^2+b^2=q\]
 \[-4a^2+a^2+b^2=q\]
 \[-3a^2+b^2=q\]
 \[b^2-3a^2=q\]Our answer is clearly $\boxed{\textbf{(E) }b^2-3a^2}$



	\item (\textit{AHSME 1972}) 
                    \begin{figure}[H]
                        \centering
                        \includegraphics[width=0.5\textwidth]{Problem/AHSME_I_1972_23_Problem_Figure_1.png}
                    \end{figure}
                    


 The radius of the smallest circle containing the symmetric figure composed of the 3 unit squares shown above is


 $\textbf{(A) }\sqrt{2}\qquad \textbf{(B) }\sqrt{1.25}\qquad \textbf{(C) }1.25\qquad \textbf{(D) }\frac{5\sqrt{17}}{16}\qquad  \textbf{(E) }\text{None of these}$

 \textbf{Answer: }D



\textbf{Solution 1}


                    \begin{figure}[H]
                        \centering
                        \includegraphics[width=0.5\textwidth]{Solution/AHSME_I_1972_23_Solution_Figure_1.png}
                    \end{figure}
                    

 Draw lines from points $A$ and $G$ to the center of the circle $K$. As shown from the diagram, $r$ is the radius of the circumscribed circle. Set $\overline{EK}$ to $x$, so as $\overline{EF}$ is $1+1=2$, $\overline{FK}$ is $2-x$. Then, $\triangle FGK$ and $\triangle AEK$ are both right, so an equation can be formed with the Pythagorean Theorem. The radii (hypotenuses) are equal, so the following equation can be made and solved:
 \begin{align*} \frac{1}{2}^2+(2-x)^2&=1^2+x^2\\ \frac{1}{4}+(-x+2)^2&=1+x^2\\ \frac{1}{4}+x^2-4x+4&=1+x^2\\ \frac{1}{4}-4x+4&=1\\ -4x&=-3-\frac{1}{4}\\ 4x&=\frac{13}{4}\\ x&=\frac{13}{16}\\ \end{align*}
$x$ is solved, so plugging it in to the original formula yields:
 \begin{align*} \sqrt{1^2+\frac{13}{16}^2} \\ =\sqrt{1+\frac{169}{256}} \\ =\sqrt{\frac{425}{256}} \\ =\frac{\sqrt{425}}{\sqrt{256}} \\ =\frac{5\sqrt{17}}{16} \\ =\boxed{\textbf{(D)}}. \end{align*}
~airbus-a321

 


	\item (\textit{AHSME 1972}) A man walked a certain distance at a constant rate. If he had gone $\textstyle\frac{1}{2}$ mile per hour faster, he would have walked the distance in four-fifths of the time; if he had gone $\textstyle\frac{1}{2}$ mile per hour slower, he would have been $2\textstyle\frac{1}{2}$ hours longer on the road. The distance in miles he walked was


 $\textbf{(A) }13\textstyle\frac{1}{2}\qquad \textbf{(B) }15\qquad \textbf{(C) }17\frac{1}{2}\qquad \textbf{(D) }20\qquad  \textbf{(E) }25$

 \textbf{Answer: }B



\textbf{Solution}

We can make three equations out of the information, and since the distances are the same, we can equate these equations.

 
 \[\frac{4t}{5}(x+\frac{1}{2})=xt=(t+\frac{5}{2})(x-\frac{1}{2})\]where $x$ is the man's rate and $t$ is the time it takes him.  

 Looking at the first two parts of the equations, 

 
 \[\frac{4t}{5}(x+\frac{1}{2})=xt\]

 we note that we can solve for $x$. Solving for $x$, we get $x=2.$

 Now we look at the last two parts of the equation:

 
 \[xt=(t+\frac{5}{2})(x-\frac{1}{2})\]

 we note that we can solve for $t$ and we get $t=\frac{15}{2}.$  We want the find the distance, which is $xt= \boxed{15}.$

 -edited for readability

 


	\item (\textit{AHSME 1972}) Inscribed in a circle is a quadrilateral having sides of lengths $25,~39,~52$, and $60$ taken consecutively. The diameter of this circle has length


 $\textbf{(A) }62\qquad \textbf{(B) }63\qquad \textbf{(C) }65\qquad \textbf{(D) }66\qquad  \textbf{(E) }69$

 \textbf{Answer: }C



\textbf{Solution}

We note that $25^2+60^2=65^2$ and $39^2+52^2=65^2$ so our answer is $\boxed{C}$.

 -Pleaseletmewin

 Alternate Solution:

 Let's call $\overline{AB}=25$, $\overline{BC}=39$, $\overline{CD}=52$, $\overline{DA}=60$. Let's call $\overline{BD}=x$ and $\angle{DAB}=y$. By LoC we get the relations

 
 \[x^2=25^2+60^2-3000\cos(y)\]
 \[x^2=39^2+52^2-4056\cos(180-y)\] 

 If we do a bit of computation we get $x^2=4225-3000\cos(y)$, and $x^2=4225-4056\cos(y)$. This means that $4056\cos(y)=3000\cos(180-y)$. 

 We know that $\cos(180-y)=-\cos(y)$ so substituting back in we get $4056\cos(y)=-3000\cos(y)$. We can clearly see that the only solution of this is $\cos(y)=0$ or $y=90$. This then means that $\angle{BAD}=90$ and $\angle{BCD}=90$. If a triangle is a right triangle and is inscribed in a circle then the diameter is the hypotenuse.

 This means that the diameter is $\sqrt{4225}=65$ so our answer is $\boxed{C}$.

 -Bole (edited for easier readability)

 


	\item (\textit{AHSME 1972}) 
                    \begin{figure}[H]
                        \centering
                        \includegraphics[width=0.5\textwidth]{Problem/AHSME_I_1972_26_Problem_Figure_1.png}
                    \end{figure}
                    


 In the circle above, $M$ is the midpoint of arc $CAB$ and segment $MP$ is perpendicular to chord $AB$ at $P$. 
If the measure of chord $AC$ is $x$ and that of segment $AP$ is $(x+1)$, then segment $PB$ has measure equal to


 $\textbf{(A) }3x+2\qquad \textbf{(B) }3x+1\qquad \textbf{(C) }2x+3\qquad \textbf{(D) }2x+2\qquad  \textbf{(E) }2x+1$

 \textbf{Answer: }E



\textbf{Solution}

$\fbox{E}$



	\item (\textit{AHSME 1972}) If the area of $\triangle ABC$ is $64$ square units and the geometric mean (mean proportional) 
between sides $AB$ and $AC$ is $12$ inches, then $\sin A$ is equal to


 $\textbf{(A) }\dfrac{\sqrt{3}}{2}\qquad \textbf{(B) }\frac{3}{5}\qquad \textbf{(C) }\frac{4}{5}\qquad \textbf{(D) }\frac{8}{9}\qquad \textbf{(E) }\frac{15}{17}$

 \textbf{Answer: }D



\textbf{Solution}

Draw Diagram later

 We can let $AB=s$ and $AC=r$. We can also say that the area of a triangle is $\frac{1}{2}rs\sin A$, which we know, is $64$. We also know that the geometric mean of $r$ and $s$ is 12, so $\sqrt{rs}=12$.

 
Squaring both sides gives us that $rs=144$. We can substitute this value into the equation we had earlier. This gives us 
 \[\frac{1}{2}\times144\times\sin A=64\]
 \[72\sin A=64\]
 \[\sin A=\frac{64}{72}=\frac{8}{9} \Rightarrow \boxed{\text{D}}\]

 -edited for readability

 


	\item (\textit{AHSME 1972}) A circular disc with diameter $D$ is placed on an $8\times 8$ checkerboard with width $D$ so that the centers coincide. The number of checkerboard squares which are completely covered by the disc is


 $\textbf{(A) }48\qquad \textbf{(B) }44\qquad \textbf{(C) }40\qquad \textbf{(D) }36\qquad  \textbf{(E) }32$

 \textbf{Answer: }E



\textbf{Solution}

Consider the upper right half of the grid, which consists of a $4\times4$ section of the checkerboard and a quarter-circle of radius $4$. We can draw this as a coordinate grid and shade in the complete squares. There are $8$ squares in the upper right corner, so there are $8 \cdot 4 = \boxed{32}$ whole squares in total.

 The answer is $\textbf{(E)}.$

 -edited by coolmath34

 


	\item (\textit{AHSME 1972}) If $f(x)=\log \left(\frac{1+x}{1-x}\right)$ for $-1<x<1$, then $f\left(\frac{3x+x^3}{1+3x^2}\right)$ in terms of $f(x)$ is


 $\textbf{(A) }-f(x)\qquad \textbf{(B) }2f(x)\qquad \textbf{(C) }3f(x)\qquad \textbf{(D) }\left[f(x)\right]^2\qquad \\ \textbf{(E) }[f(x)]^3-f(x)$

 \textbf{Answer: }C



\textbf{Solution}

For simplicity, we can write $f(x) = \log(\frac{1+x}{1-x})$ as $f(x) = \log({1+x}) - \log({1-x})$ via the divisibility log property. From here we can plug in $\frac{3x+x^3}{1+3x^2}$ for $x$ which becomes:

 $\log({1+\frac{3x+x^3}{1+3x^2})} - \log({1-\frac{3x+x^3}{1+3x^2})} \implies$
$\log{(\frac{1+3x^2}{1+3x^2}+\frac{3x+x^3}{1+3x^2})} - \log{(\frac{1+3x^2}{1+3x^2}-\frac{3x+x^3}{1+3x^2})} \implies$
$\log{(\frac{x^3+3x^2+3x+1}{1+3x^2})} - \log{(\frac{-x^3+3x^2-3x+1}{1+3x^2})} \implies$
$\log{(\frac{(x+1)^3}{1+3x^2})} - \log{(\frac{(-x+1)^3}{1+3x^2})}$

 Now, if we bring the two "$\log$" terms back together ($f(x) = \log\left(\dfrac{\dfrac{(x+1)^3}{3x^2+1}}{\dfrac{(-x+1)^3} {3x^2+1}}\right)$), it becomes apparent that the two denominators cancel out, making for:

 $f(x) = \log(\frac{(x+1)^3}{(-x+1)^3}) \implies$
$f(x) = 3\log(\frac{1+x}{1-x})$ via the power law. We can then identify that $\frac{1+x}{1-x}$ is $f(x)$, and by substitution, the answer is $3f(x)$ or $\boxed{\textbf{(C)}}$



	\item (\textit{AHSME 1972}) 
                    \begin{figure}[H]
                        \centering
                        \includegraphics[width=0.5\textwidth]{Problem/AHSME_I_1972_30_Problem_Figure_1.png}
                    \end{figure}
                    


 A rectangular piece of paper 6 inches wide is folded as in the diagram so that one corner touches the opposite side. The length in inches of the crease L in terms of angle $\theta$ is


 $\textbf{(A) }3\sec ^2\theta\csc\theta\qquad \textbf{(B) }6\sin\theta\sec\theta\qquad \textbf{(C) }3\sec\theta\csc\theta\qquad \textbf{(D) }6\sec\theta\csc^2\theta\qquad \textbf{(E) }\text{None of these}$

 \textbf{Answer: }A



\textbf{Solution}


                    \begin{figure}[H]
                        \centering
                        \includegraphics[width=0.5\textwidth]{Solution/AHSME_I_1972_30_Solution_Figure_1.png}
                    \end{figure}
                    Let the rectangle $ABCD$ have crease $BE$ with $E$ on $CD$, and let $F$ be on $AD$ such that $F$ is a reflection of $C$ over $BE$.  Notice that triangles $ABF$ and $DEF$ are similar, so by setting $CE = EF = x$ with $DE = 6-x$, giving $DF = 2\sqrt{3x-9}$ we have that $AF = \frac{18-3x}{\sqrt{3x-9}}$. Noticing that $BC = BF = x\cot{\theta}$ gives $\frac{(18-3x)^2}{3x-9}+36 = x^2\cot^2{\theta} \Rightarrow \frac{3(x-6)^2}{x-3}+36 = \frac{3x^2}{x-3} = x^2\cot^2{\theta}$. $x = \frac{3\cot^2{\theta}+3}{\cot^2{\theta}} = \frac{3\csc^2{\theta}}{\cot^2{\theta}} = 3\sec^2{\theta}$. Noticing that $BE = x\sqrt{\cot^2{\theta}+1} = x\csc{\theta}$ gives the answer to be $3\sec^2{\theta}\csc{\theta}$.

 Thanks to xiej for correcting my mistake!

 


	\item (\textit{AHSME 1972}) When the number $2^{1000}$ is divided by $13$, the remainder in the division is


 $\textbf{(A) }1\qquad \textbf{(B) }2\qquad \textbf{(C) }3\qquad \textbf{(D) }7\qquad  \textbf{(E) }11$

 \textbf{Answer: }C



\textbf{Solution}

By \href{https://artofproblemsolving.com/wiki/index.php?title=Fermat\%27s\_Little\_Theorem}{Fermat's Little Theorem}, we know that $2^{1000} \equiv 2^{1000 \pmod{12}}\pmod{13}$. However, we find that $1000 \equiv 4 \pmod{12}$, so $2^{1000} \equiv 2^4 = 16 \equiv 3 \pmod{13}$, so the answer is $\boxed{\textbf{(C)}}$.

 


	\item (\textit{AHSME 1972}) 
                    \begin{figure}[H]
                        \centering
                        \includegraphics[width=0.5\textwidth]{Problem/AHSME_I_1972_32_Problem_Figure_1.png}
                    \end{figure}
                    


 Chords $AB$ and $CD$ in the circle above intersect at E and are perpendicular to each other. 
If segments $AE, EB$, and $ED$ have measures $2, 3$, and $6$ respectively, then the length of the diameter of the circle is


 $\textbf{(A) }4\sqrt{5}\qquad \textbf{(B) }\sqrt{65}\qquad \textbf{(C) }2\sqrt{17}\qquad \textbf{(D) }3\sqrt{7}\qquad  \textbf{(E) }6\sqrt{2}$

 \textbf{Answer: }B



\textbf{Solution}

Using the chord theorem we can immediately solve for $\overline {CE}$ which will help us. The chord theorem states that if two chords, $\overline {AB}$ and $\overline {CD}$, intersect at lets say $E$, then $\overline {AE} \cdot \overline {EB}=\overline {CE} \cdot \overline {ED}$. Now that we know this we can label length $\overline {CE}$ as $\eta$. We now have the equation $3\eta=2\cdot6\implies3\eta=12\implies\eta=4$. We now know the chord lengths $7$ and $8$. We now label the center of the circle as $O$. If the circle was on the graph, with $O$ being right on the origin, then make a $x=0$ line cutting through the center. This bisects chord $AB$ at point $M$. This means $\overline {AM}=4$ and $\overline {EM}=2$ If we now draw a "$y=0$" line which bisects chord $CD$ at point $N$ making $\overline {CN}=\frac{7}{2}$. If we now move up $\overline {EM}$ to $\overline {ON}$ then we can find the radius($\overline {CO}$) using Pythagorean Theorem.

 $\sqrt{{\overline {EM}}^2+{\overline {ON}}^2}=r$
$\sqrt{2^2+(\frac{7}{2})^2}$
$\sqrt{4+\frac{49}{4}}$
$\sqrt{\frac{65}{4}}$
$\frac{\sqrt{65}}{2}=r$

 Now that we have $r$, the diameter is just two times of that, so $\boxed{\textbf{(B)}\sqrt{65}}$



	\item (\textit{AHSME 1972}) The minimum value of the quotient of a (base ten) number of three different non-zero digits divided by the sum of its digits is


 $\textbf{(A) }9.7\qquad \textbf{(B) }10.1\qquad \textbf{(C) }10.5\qquad \textbf{(D) }10.9\qquad  \textbf{(E) }20.5$

 \textbf{Answer: }C



\textbf{Solution}

The answer we are looking for can be expressed as $\dfrac{100a+10b+c}{a+b+c}$. This is equivalent to $1 + \dfrac{99a+9b}{a+b+c}$. Because we are trying to minimize our solution, we set $c$ = $9$, so we have $1 + \dfrac{99a+9b}{a+b+9}$. This is equal to $1 + \dfrac{9a+9b+81}{a+b+9} + \dfrac{90a-81}{a+b+9}$, which simplifies to $10+ \dfrac{90a-81}{a+b+9}$. Since each digit is unique, we set $b$ to $8$, leaving us with $10 + \dfrac{90a-81}{a+17}$. Clearly, $a$ should be minimized, so $a = 1$ and our answer is 
 \[\boxed{\textbf{(C) }10.5}.\]



	\item (\textit{AHSME 1972}) Three times Dick's age plus Tom's age equals twice Harry's age. 
Double the cube of Harry's age is equal to three times the cube of Dick's age added to the cube of Tom's age. 
Their respective ages are relatively prime to each other. The sum of the squares of their ages is


 $\textbf{(A) }42\qquad \textbf{(B) }46\qquad \textbf{(C) }122\qquad \textbf{(D) }290\qquad  \textbf{(E) }326$

 \textbf{Answer: }A



\textbf{Solution}


 \[t=2h-3d\]
 \[3d^3+t^3=2h^3\]

 First, substitute in t into the second equation and get $3d^3+8h^3-36h^2d+54hd^2-27d^3=2h^3$. That turns into $h^3-6h^2d+9hd^2-4d^3=0$ which is factored into $(h-4d)(h-d)^2 =0.$ WLOG, $d=1$ and consequently $h=4$. Then $t=8-3=5$. Everything appears to be relatively prime already. The answer is thus $1+16+25=\boxed{\textbf{(A) }42}.$ ~lopkiloinm

 


	\item (\textit{AHSME 1972}) 
                    \begin{figure}[H]
                        \centering
                        \includegraphics[width=0.5\textwidth]{Problem/AHSME_I_1972_35_Problem_Figure_1.png}
                    \end{figure}
                    


 Equilateral triangle $ABP$ (see figure) with side $AB$ of length $2$ inches is placed inside square $AXYZ$ with side of length $4$ 
inches so that $B$ is on side $AX$. The triangle is rotated clockwise about $B$, then $P$, and so on along the sides of the square 
until $P$ returns to its original position. The length of the path in inches traversed by vertex $P$ is equal to


 $\textbf{(A) }20\pi/3\qquad \textbf{(B) }32\pi/3\qquad \textbf{(C) }12\pi\qquad \textbf{(D) }40\pi/3\qquad  \textbf{(E) }15\pi$

 \textbf{Answer: }D



\textbf{Solution}

$\fbox{D}$



\end{enumerate}